\subsection{Evaluation protocol}
\label{sec:experiments}

The model interface uses OpenAI Codex CLI \citep{openai-codex} through a custom
episode harness. Software version evidence is documented in
Appendix~\ref{app:software-versions}.

% EDITING NOTE (retained): The pipeline -----should have a hand drawn graph
\subsubsection{Episode procedure}
\label{sec:exp-pipeline}

One episode is one sample attempted by one model under one tier of assistance. The loop is the same
for every arm and every model, and the only thing that varies between arms is what the model is
allowed to see and call.

Each query contains one tool call. Accepted edits update the current sequence;
rejected edits preserve it, and both consume budget. Removing a fold replays the remaining actions;
returning to an earlier step truncates the sequence, while revision restoration can recover a
previous branch. The reference sequence is withheld until scoring. Run-specific turn, time, and
cycle limits must be reported alongside the query count.

% ===========================================================================
% FIGURE 4 -- THE EVALUATION PIPELINE.  [TO DRAW]
% Left-to-right diagram of one episode.
%   LEFT    the two inputs: crease pattern, final folded state.
%   CENTRE  the loop as a cycle: model proposes a fold; environment returns either
%           a new state with its two rendered views, or a NAMED refusal; history
%           accumulates.
%   RIGHT   termination and scoring: replay of the model's sequence and replay of
%           the reference, compared at Level 1 and Level 2.
% Must make three things visually obvious:
%   (1) environment and verifier are ONE box, not two;
%   (2) refusals loop back WITHOUT advancing;
%   (3) seq.json sits OUTSIDE the loop, entering only at scoring.
% A dashed boundary around the withheld items (intermediate frames, reference
% sequence) carries point (3) without a sentence.
% ===========================================================================
Episode pseudocode and the evaluation pipeline appear in
Appendix~\ref{app:evaluation-pipeline}.

\subsubsection{Tool configuration}
\label{sec:exp-tools}

Table~\ref{tab:tools} lists the tool interface. Basic tools comprise all entries except
the optional \texttt{list\_legal\_folds} and \texttt{compare\_to\_target} tools.
Loading and configuration are described in
Appendix~\ref{app:tool-loading}.

\begin{table}[h]
\centering
\caption{Tools available to the evaluation harness; exposure depends on run configuration.}
\label{tab:tools}
\begin{tabular}{lp{0.65\textwidth}}
\hline
Tool & Operation \\
\hline
\texttt{add\_fold} & Append a checked fold; reject invalid actions without changing state. \\
\texttt{remove\_fold} & Remove a specified fold and replay the remainder transactionally. \\
\texttt{go\_to\_step} & Retain the first specified number of folds. \\
\texttt{restore\_revision} & Restore a previously recorded branch state. \\
\texttt{get\_state} & Return the current sequence and layer geometry. \\
\texttt{get\_images} & Render top, oblique, reverse, X-ray, and exploded views of a completed step; exposed views depend on run configuration. \\
\texttt{list\_legal\_folds} & Optionally enumerate legal candidates compatible with the target CP. \\
\texttt{compare\_to\_target} & Optionally return feedback at the configured level (Section~\ref{sec:target-feedback}). \\
\texttt{finish} & End the episode and evaluate the current sequence. \\
\hline
\end{tabular}
\end{table}

% ===========================================================================
% FIGURE 5 -- A REFUSAL AS THE MODEL RECEIVES IT.  [TO DRAW]
% One concrete rejected fold, rendered exactly as the harness returns it:
%   - the proposed fold line drawn over the current state;
%   - the named reason (would-tear);
%   - the specific evidence: for a tear, the crease segment joining a moving face
%     to a stationary one, highlighted, away from the fold line.
% Beside it, the SAME state with a legal fold line, for contrast. This is the
% tearing pair of the action-space section, which currently exists only as prose.
% Caption states that the reason is a NAME rather than a boolean, and that this is
% why error analysis is a taxonomy that exists before the experiment.
% ===========================================================================
% EDITING NOTE: Refusal figure moved to the appendix plan; use an actual rejected action.


\subsubsection{Models and budgets}
\label{sec:exp-protocol}

The recent ablations evaluate GPT-5.6 Luna at low reasoning effort and GPT-6 Luna
at low and high effort, using ten samples per group and configuration. Runs use an
80-turn limit, a 300-second call timeout, and full image history. Earlier runs have
different coverage and budgets and are reported separately (Section~\ref{sec:results}).
% EDITING NOTE: Report per-run stopping limits and repeat variability from the inventory.

We include breadth-first search over the same enumerated legal actions as a
systematic-search baseline (Appendix~\ref{app:bfs-analysis}).
% EDITING NOTE: Planned baselines retained for future work: uniform random search
% over legal candidates and the symbolic crease-pattern solver cited as creasy.
